% -----------------------------------------------------------------------------
% One-loop QED: pole, charge and vertex renormalization as one physical chain
% -----------------------------------------------------------------------------

树级 QED 只有一个电子--光子局域顶角，但一圈以后出现三个彼此不同的一粒子不可约基本核：电子自能 \(\Sigma\)、光子真空极化 \(\Pi^{\mu\nu}\) 与电磁顶角修正 \(\Lambda^\mu\)。它们分别对应三类可测信息：电子传播子的单粒子极点确定物理质量与外腿归一化；长波长电磁交换确定低能电荷；在壳电磁顶角的额外张量结构确定磁矩等自旋响应。紫外发散出现在裸参数与这些可测定义之间的微扰关系中，质量、电荷和场强反项由相应的极点、留数与低能归一化条件固定。

对以特征物理动量 \(Q\) 探测、且没有大对数重求和的过程，固定一圈截断的两个控制量为
\begin{equation*}
\epsilon_{\rm 1L}\equiv\frac{\alpha(\mu)}{\pi},
\qquad
\epsilon_{\log}\equiv
\frac{\alpha(\mu)}{\pi}
\left|\ln\frac{Q^2}{\me^2c^2}\right|.
\end{equation*}
当 \(\epsilon_{\rm 1L}\ll1\) 且 \(\epsilon_{\log}\ll1\) 时，一圈截断把普通微扰误差控制在下一圈阶；若 \(\epsilon_{\log}\gtrsim1\)，固定圈阶中的大对数必须用重整化群或相应因子化结构重求和。Ward--Takahashi 恒等式本身不依赖这一截断。

电子二点函数先给出质量与外腿归一化。把重整化后的约化逆传播子写成
\begin{equation}
\boxed{
S_R^{-1}(p)
=\slashed p-\me c-\Sigma_R(p).}
\label{eq:qed-renormalized-electron-inverse-dp53}
\end{equation}
一般 Lorentz 协变自能同时含 $\slashed p$ 与标量两种 Dirac 结构，因此记为 $\Sigma_R(p)$。为定义单粒子极点沿质量壳法向的斜率，固定任意未来指向单位类时四矢量 $u^\mu$，满足 $u^2=1$，并沿同一方向写
\begin{equation*}
 p^\mu=\lambda u^\mu,
 \qquad [\lambda]={\rm kg\,m\,s^{-1}}.
\end{equation*}
正能投影算符为
\begin{equation*}
 P_+(u)\equiv\frac{\id_4+\slashed u}{2},
 \qquad P_+^2=P_+.
\end{equation*}
Lorentz 协变性保证自能在这个二维正能子空间内只留下一个标量本征值；定义
\begin{equation*}
 \sigma_R(\lambda)
 \equiv\frac12\Tr\!\left[P_+(u)\Sigma_R(\lambda u)\right],
 \qquad
 P_+\Sigma_R(\lambda u)P_+=\sigma_R(\lambda)P_+.
\end{equation*}
在壳点 $\lambda=\me c$，物理质量要求投影逆传播子在那里有零点，而单位留数要求这个零点对标量极点坐标 $\lambda$ 的一阶斜率等于自由值。两条条件严格写成
\begin{equation}
\boxed{
\sigma_R(\me c)=0,
\qquad
\left.\frac{\dd\sigma_R(\lambda)}{\dd\lambda}\right|_{\lambda=\me c}=0.}
\label{eq:qed-onshell-selfenergy-conditions-dp53}
\end{equation}
第一式固定质量反项，使极点仍在 $p^2=\me^2c^2$；第二式固定电子场强反项，使 LSZ 外腿抽出的正能单粒子极点具有单位留数。传统写法
$\Sigma_R(\slashed p)|_{\slashed p=\me c}=0$ 与
$\partial\Sigma_R/\partial\slashed p|_{\me c}=0$
就是这两个投影标量条件的简写；所谓“对 $\slashed p$ 求导”因此始终指这两个 Lorentz 标量系数的普通导数，而不是未定义的矩阵微分。

光子二点函数则固定电荷的低能归一化。真空极化的横向分解为
\begin{equation*}
\Pi^{\mu\nu}(q)
=(q^2\eta^{\mu\nu}-q^\mu q^\nu)\Pi(q^2),
\end{equation*}
在壳电荷方案中取
\begin{equation}
\boxed{\Pi_R(0)=0.}
\label{eq:qed-onshell-vacpol-condition-dp53}
\end{equation}
于是 \(q^2\to0\) 的长程 Coulomb/Thomson 极限直接使用实验测得的低能电荷；非零 \(q^2\) 时剩余的 \(\Pi_R(q^2)\) 才描述真空屏蔽随分辨尺度改变的物理效应。

电磁顶角把前两个问题连接起来。对在壳电子外腿，完整顶角由两个无量纲形状因子参数化：
\begin{equation*}
\bar u(p')\Gamma_R^\mu(p',p)u(p)
=\bar u(p')\left[
F_1(q^2)\gamma^\mu
+\frac{\ii\sigma^{\mu\nu}q_\nu}{2\me c}F_2(q^2)
\right]u(p).
\end{equation*}
低能电荷的定义要求
\begin{equation}
\boxed{F_1(0)=1,}
\label{eq:qed-F1-charge-normalization-dp53}
\end{equation}
而 \(F_2(0)\) 是 Ward--Takahashi 恒等式没有固定的有限横向结构；它给出反常磁矩
\begin{equation*}
a_e=\frac{g_{\rm mag}-2}{2}=F_2(0).
\end{equation*}
因此质量、低能电荷和反常磁矩分别由电子极点、零动量光子交换和横向顶角结构定义，对应传播子和顶角的三个独立可测信息。

精确 Ward--Takahashi 恒等\eqnrefs{eq:qed_ward_takahashi_momentum} 在一圈阶把电子自能与零动量转移顶角核直接联系起来。两者写为
\begin{equation}
\boxed{
\begin{aligned}
\Sigma_{\rm 1L}(p)
&=-\ii g_{\rm em}^2
\int_{\ell}^{(d)}
\frac{\gamma^\alpha\mathsf S_0(p-\ell)\gamma_\alpha}
{\ell^2+\ii0^+},\\
\Lambda_{\rm 1L}^{\mu}(p,p)
&=-\ii g_{\rm em}^2
\int_{\ell}^{(d)}
\frac{\gamma^\alpha\mathsf S_0(p-\ell)\gamma^\mu
\mathsf S_0(p-\ell)\gamma_\alpha}
{\ell^2+\ii0^+},
\end{aligned}}
\label{eq:qed-one-loop-selfenergy-vertex-kernels-dp68}
\end{equation}
其中 \(\mathsf S_0(k)\) 采用自由电子约化传播核及其既定极点处方。两式的区别只有内部电子线上是否插入一个零动量光子顶角。对同一内部传播核关于外动量求导，正好产生这个插入，因此
\begin{equation}
\boxed{
\Lambda_{\rm 1L}^{\mu}(p,p)
=-\frac{\partial\Sigma_{\rm 1L}(p)}{\partial p_\mu}.}
\label{eq:qed-one-loop-ward-kernel-dp53}
\end{equation}
这是一圈 Ward--Takahashi 恒等式的图级形式：电子自能对外动量的变化与零动量电磁顶角不是两个独立的紫外结构。

自能中乘在 \(\slashed p\) 前的紫外发散因此必须与最小耦合顶角的纵向紫外发散具有同一系数。若 \(Z_2\) 与 \(Z_1\) 分别表示电子场和顶角的重整化常数，则
\begin{equation}
\boxed{Z_1=Z_2.}
\label{eq:qed-Z1-Z2-dp53}
\end{equation}

电荷关系还需要比较同一个裸相互作用在裸场与重整化场中的两种写法。由
\(\psi_0=Z_2^{1/2}\psi_R\)、\(A_{0\mu}=Z_3^{1/2}A_{R\mu}\)，裸最小耦合项与以 \(Z_1\) 参数化的重整化顶角具有相同系数，当且仅当
\begin{equation}
\boxed{
q_0 Z_2 Z_3^{1/2}=q_R Z_1.}
\label{eq:qed-bare-renormalized-charge-match-dp68}
\end{equation}
代入\eqnrefs{eq:qed-Z1-Z2-dp53} 后，电子场因子约去，得到
\begin{equation}
\boxed{q_0=q_RZ_3^{-1/2}.}
\label{eq:qed-charge-Z3-dp53}
\end{equation}
所以阿贝尔电荷重整化的独立紫外信息只剩光子场因子 \(Z_3\)：电子外腿与纵向顶角的紫外重整化已经被局域 \(U(1)\) 对称性锁定为同一个量。

若把一圈 QED 看成对树级理论的第一次形变，可以用三个低能实验量给每个基本核定位。电子静止能确定精确传播子的极点位置 \(E=\me c^2\)，因此电子自能的常数部分不能独立成为新参数，而被质量反项重新吸收到物理质量定义中。低动量 Coulomb 或 Thomson 散射确定 \(q_R\)，因此真空极化在 \(q^2=0\) 的常数部分被条件 \(\Pi_R(0)=0\) 吸收；有限 \(Q\) 的偏离才表现为运行耦合。电子自旋在弱磁场中的进动则测量 \(g_{\rm mag}\)，其树级值 \(2\) 由 Dirac 方程给出，而一圈顶角的有限 Pauli 部分给出
\[
F_2(0)=\frac{\alpha}{2\pi}+O(\alpha^2).
\]
这三个观测量分别钉住极点、低能电荷和横向磁矩结构，因此一圈重整化的物理输入与三个一粒子不可约基本核的职责一一对应。
